% ----------------------------------------------------------
% Introdução
% ----------------------------------------------------------
\chapter{Introdução}
\label{cap:introducao}
% ----------------------------------------------------------
A teoria da computação estuda as propriedades dos sistemas computacionais. Essa teoria foi formalizada de forma efetiva, pela primeira vez, em 1936, por Alan Turing e Alonzo Church, que construíram as bases do que hoje chamamos de computação por meio da elaboração do conceito de linguagens formais.

Linguagens formais são reconhecidas por meio de máquinas de estados e descritas por meio de gramáticas formais. Neste trabalho, serão enfatizadas as máquinas de estados, mais especificamente os Autômatos de Pilha Não Determinísticos ou \ac{APN}, por apresentarem maior nível de abstração e representação mais complexa do que as demais máquinas.

Sendo a teoria da computação a base das tecnologias atuais, torna-se essencial que os discentes dos cursos da área de computação, como Sistemas de Informação e Engenharia da Computação, dominem esses conceitos para se tornarem profissionais qualificados. Afinal, tais conceitos não apenas fornecem o embasamento teórico necessário, mas também promovem habilidades de pensamento crítico e análise rigorosa, conferindo ao estudante a capacidade de formalizar problemas, definir algoritmos e analisar sua eficiência.

Como forma de auxiliar os discentes a visualizarem de maneira mais clara os autômatos de pilha não determinísticos, este trabalho tem como objetivo desenvolver uma ferramenta web para a modelagem e visualização de APNs em execução.

% ----------------------------

\section{Organização do capítulo}

Neste capítulo, será abordada a motivação deste trabalho na Seção \ref{sec:problema}, os objetivos propostos na Seção \ref{sec:objetivos} e a metodologia adotada para o desenvolvimento do projeto na Seção \ref{sec:metodologia}.

\section{O problema de pesquisa}
\label{sec:problema}

Ciente da importância do estudo de linguagens formais e máquinas de estados — utilizadas hoje como ferramentas fundamentais para o desenvolvimento de compiladores de linguagens de programação modernas —, é imprescindível que os alunos dos cursos superiores da área de computação dominem esses assuntos. Contudo, eles normalmente enfrentam desafios complexos durante esse aprendizado.

Linguagens formais e máquinas de estados são conceitos teóricos rígidos e abstratos. Mesmo com o auxílio de diagramas, o processo de ensino-aprendizagem costuma ser desafiador \cite{Nelishia2010}, especialmente no que tange às linguagens livres de contexto ou \ac{LLC}. A máquina reconhecedora dessas linguagens, o Autômato de Pilha Não Determinístico, possui uma série de execuções simultâneas em que cada ramificação carrega seu próprio conjunto de informações (como o topo da pilha e o estado atual). Assimilar toda essa execução paralela de forma puramente teórica constitui um grande obstáculo no aprendizado da disciplina.

\section{Objetivos}
\label{sec:objetivos}

Dadas as dificuldades dos alunos na assimilação das linguagens livres de contexto e de suas respectivas máquinas de estados, este trabalho visa criar uma plataforma web na qual seja possível projetar esses modelos e demonstrar sua execução em tempo real. Desse modo, busca-se tornar viável a visualização detalhada de cada uma das múltiplas execuções simultâneas dos autômatos de pilha não determinísticos.

Para alcançar o objetivo geral, este trabalho possui os seguintes objetivos específicos:

\begin{itemize}
    \item Criar um software gráfico que auxilie os alunos a visualizar a execução de autômatos de pilha em tempo real;
    \item Explorar as derivadas propostas por \cite{Brzozowski64} para a geração de palavras pertencentes a uma linguagem formal;
    \item Disponibilizar o sistema de forma web e gratuita para facilitar ao máximo o seu acesso. 
\end{itemize}

\section{Metodologia}
\label{sec:metodologia}

A metodologia deste trabalho consiste no desenvolvimento de um sistema de criação e visualização de autômatos e de sua respectiva execução, servindo como um meio didático para que o estudante consiga identificar erros em suas tarefas de forma rápida, visual e automatizada.

\section{Organização do trabalho}

O Capítulo~\ref{cap:revisao} apresenta os conceitos básicos necessários para a compreensão do trabalho, bem como as ferramentas já existentes com fins semelhantes; o Capítulo~\ref{cap:desenvolvimento} detalha as ideias gerais da implementação da ferramenta de simulação proposta; o Capítulo~\ref{cap:resultados} discute brevemente os resultados obtidos durante o desenvolvimento; por fim, o Capítulo~\ref{cap:conclusao} apresenta as conclusões e as considerações finais deste trabalho.
